% Pending method-obstruction note; not an included result or A5 gate movement.
\subsection*{Properness, openness and pure two-dimensional fibres do not ensure a Thom stratification}

Let \(K\) be a real closed field, write \(C=K(i)\), and regard
\(C\) as \(K^2\).  Consider
\[
 Y=C\times K,\qquad
 X=\{(a,b,[u:v:w])\in Y\times\mathbb P^2(C):
                   u^2-a^2vw-bw^2=0\},
 \qquad \pi:X\longrightarrow Y.
 \tag{T.1}
\]
Here complex projective space is viewed as a compact semialgebraic
space, for example through its realization as the rank-one Hermitian
orthogonal projection matrices of trace one.  The homogeneous equation
defines a closed semialgebraic subset of the product.  The map \(\pi\)
is surjective, since \([0:1:0]\) is in every fibre, and definably proper,
since the second factor is definably compact.

Every fibre has pure real dimension two.  When \(a\ne0\), it is a smooth
projective conic.  When \(a=0\), \(b\ne0\), it is the union of the
two distinct complex projective lines
\(u=\sqrt b\,w\) and \(u=-\sqrt b\,w\), meeting at \([0:1:0]\).
When \(a=b=0\), its underlying point set is the projective line \(u=0\).
All square roots in this paragraph are taken in the algebraically
closed field \(C\); no choice of a continuous square root is required.

The map \(\pi\) is open.  In the chart \(w=1\), put \(z=v\), so the
equation is \(u^2=a^2z+b\).  At a point with coordinates
\((a_0,b_0,u_0,z_0)\), fix \(z=z_0\).  For all \((a,b)\) sufficiently
close to \((a_0,b_0)\), there is a square root of \(a^2z_0+b\)
arbitrarily close to \(u_0\).  If \(u_0\ne0\), this is the elementary
local inverse theorem for the nonsingular map \(u\mapsto u^2\)
of \(C\); if \(u_0=0\), every square root has modulus
\(\sqrt{|a^2z_0+b|}\), which tends to zero.  Thus the image of every
neighbourhood of that point contains a target neighbourhood.
The only remaining projective point is \([0:1:0]\); there openness
follows from the continuous constant projective section
\((a,b)\mapsto(a,b,[0:1:0])\) through the point.

Nevertheless, \(\pi\) admits no finite definable \(C^1\) Thom
\(a_\pi\) stratification.  In the affine chart above, its central
fibre is
\[
 E=\{a=b=u=0,\ z\in C\}.
\]
In any hypothetical finite stratification, one source stratum \(T\)
has an intersection with \(E\) containing a nonempty relatively open
subset of \(E\).  This follows by finite definable dimension: otherwise
finitely many subsets of dimension at most one would cover \(E\).
Choose \(q=(0,0,0,z_0)\) in that relative interior with \(z_0\ne0\).
The tangent plane \(C_z=T_qE\) is contained in
\(\ker d(\pi|_T)_q\), since \(\pi\) is constant on \(E\).

Choose \(\alpha\in C^*\) with \(\alpha^2z_0=1\), and approach \(q\)
by the definable curve
\[
 q_t=(t\alpha,-t^2,0,z_0),\qquad t>0.
 \tag{T.2}
\]
It lies in \(X\), and \(a\ne0\) on this curve.  At every \(q_t\),
the full source is a smooth real five-manifold and \(\pi\) is a
submersion.  Its vertical tangent space is exactly the complex
\(u\)-line \(C_u\).  Indeed, with \(a,b\) fixed, differentiation of
\(u^2-a^2z-b=0\) gives \(2u\,du-a^2dz=0\); at \(u=0,a\ne0\)
this is \(dz=0\).

By definability and finiteness, after shortening the interval the
curve \(q_t\) lies in a single source stratum \(S\).  Every
\(\ker d(\pi|_S)_{q_t}\) is a subspace of \(C_u\).  These kernels
have constant dimension, because the restriction to a source stratum
is a submersion to one target stratum, and their Grassmannian limit
\(L\) exists after taking a definable subcurve.  Hence \(L\subset C_u\).
The stratum \(S\) cannot equal \(T\): in that case continuity of the
constant-rank kernel bundle on \(T\) would give
\(\ker d(\pi|_T)_q\subset C_u\), contradicting its containment of
the nonzero, transverse plane \(C_z\).  Since \(q\in T\cap\overline S\),
the frontier condition gives \(T\subset\overline S\).
Thom's condition would then require
\[
                 C_z\subset\ker d(\pi|_T)_q
                          \subset L\subset C_u,
\]
which is impossible.  This proof does not assume in advance that
the central fibre is a union of source strata.

This is a method obstruction in precisely target dimension three
and fibre dimension two.  It rules out repairing the universal
Thom-stratification premise merely by adding openness and pure,
constant fibre dimension.  It is not a counterexample to the A.5
maximal-locus assertion: its maximal open local-product locus is
the explicitly definable set \(\{a\ne0\}\).  The trivialization there
is the projective linear change
\[
 [u:v:w]\longmapsto
       [U:V:W]=[u/a:v+(b/a^2)w:w],
 \qquad U^2=VW,
 \tag{T.3}
\]
with inverse \([U:V:W]\mapsto[aU:V-(b/a^2)W:W]\).
At \(a=0,b\ne0\), the fibre is two spheres meeting in one point;
deleting that point disconnects it, whereas a sphere minus any
point is connected.  Thus it is not homeomorphic to the nearby
smooth conic.  Every neighbourhood of \((0,0)\) likewise contains
both such reducible fibres and smooth conics, so it cannot be a
product neighbourhood either.

The affine equation is motivated by the original non-Thom example
of L\^e recorded by Sabbah, \emph{Morphismes analytiques stratifi\'es
sans \'eclatement et cycles \'evanescents}, Ast\'erisque 101--102
(1983), printed pp.286,288.  Sabbah's printed p.291 also explicitly
distinguishes flatness from the needed no-blowing-up condition and
records a proper flat counterexample of Deligne.  Those original
pages were opened at
\url{https://www.numdam.org/item/AST_1983__101-102__286_0.pdf}.
The projective real-three-dimensional-base construction (T.1),
its openness proof and its exact maximal locus are proved here,
and are not attributed to Sabbah.

The finite algebra certificate \path{code/check_open_conic_Thom.py}
checks the projective change and its inverse, the curve equation,
and the real vertical derivative determinant over \(\mathbb Q\).
Its actual output and runtime are in
\path{logs/check_open_conic_Thom.txt}.
